% Discusion - cada bloque aporta un parrafo; D lo une
% D ya puso el suyo. A, B y C agregan el de ellos arriba del de D.

\section{Discusión}
\label{sec:discusion}

\pendiente{A: un párrafo sobre la elección de $f_s$ y $N$, y qué tan bien se
sostiene la cuantización a 8 bits del emisor.}

\pendiente{B: un párrafo sobre el sistema completo en alto nivel: qué se
reutilizó de la Tarea 1 y qué hubo que cambiar.}

\pendiente{C: un párrafo sobre el emisor en la Pico: el DDS, el filtro RC medido
y las diferencias con lo que predecía el prototipo.}

\paragraph{Validación en el receptor.} La estrategia quedó definida con sus cuatro
diagramas y sus tres criterios de aceptación fijados antes de medir, que es lo que
pide la sección 3.5 del proyecto. Dos decisiones son las que más peso tienen.
La primera es comparar los niveles en relativo y no en absoluto: entre el DDS y la
FFT hay una ganancia desconocida que incluye el nivel del PWM, la caída del RC y
la escala del ADC, y normalizar cada espectro a su pico más alto vuelve la prueba
insensible a esa ganancia sin perder sensibilidad a la forma del espectro, que es
lo que de verdad se quiere validar. La segunda es mandar solo los picos en lugar
del espectro completo: baja el tráfico de 32\,kB/s a 750\,B/s sin perder nada que
la comparación use.

La limitación que queda es que todo esto está verificado contra capturas
simuladas, no contra la placa. La regla de decisión se probó generando
perturbaciones controladas --- un pico corrido tres bins, un nivel cinco
decibelios fuera --- y comprobando que el script las marca y dice por qué, pero el
comportamiento real de la cadena PWM--RC--ADC todavía no se ha medido. Lo primero
que habrá que revisar cuando haya hardware es si la SNR del modo 1 llega a los
30\,dB que se fijaron como mínimo, porque ese umbral se calculó a partir del techo
teórico de los 8 bits del emisor y es el supuesto más frágil de los tres.

Para el proyecto hay una consecuencia concreta: el filtro RC de la salida del
emisor sirve como reconstrucción pero no como antialias, porque su frecuencia de
corte está por encima de Nyquist. En este taller no importa, porque el emisor es
sintético y no genera nada entre la banda útil y la portadora, pero cuando el nodo
1 tome la señal de una fuente arbitraria va a necesitar un antialias propio antes
del ADC.
